\documentclass[10pt,a4paper]{amsart}
\usepackage[margin=19mm]{geometry}
\usepackage{amsmath,amssymb,mathtools}
\usepackage[scr=rsfs,cal=euler]{mathalfa}
\usepackage{xfrac}
\usepackage[unicode,hidelinks,bookmarks=false]{hyperref}
\newcommand{\ie}{i.e.\ }
\newcommand{\cf}{cf.\ }
\newcommand{\resp}{respectively\ }
\newcommand{\Subsection}{\subsection}
\providecommand{\nobreakdash}{\nobreak\hbox{-}}
\setlength{\emergencystretch}{2em}
\multlinegap0pt
\usepackage[shortlabels]{enumitem}
\usepackage{amssymb}
\usepackage{mathtools}
\newtheorem{lemma}{Lemma}
\newtheorem{theorem}{Theorem}
\usepackage{cleveref}
\crefname{lemma}{Lemma}{Lemmas}
\crefname{theorem}{Theorem}{Theorems}
\newcommand{\Beta}{\operatorname{Beta}}
\newcommand{\Cauchy}{\operatorname{Cauchy}}
\newcommand{\E}{\mathbb E}
\newcommand{\Law}{\mathcal L}
\newcommand{\R}{\mathbb R}
\newcommand{\dd}{\,\mathrm d}
\newcommand{\sphere}{\mathbb S}
\title{Universal Beta Incidence Angles: Cauchy Rigidity and\\ Infinite Arrangements}
\author{Tianle Liu}
\date{Source: arXiv:2609.00603v1 (CC BY 4.0)}
\begin{document}
\maketitle

\noindent\emph{Source and licence.} Universal Beta Incidence Angles: Cauchy Rigidity and\\ Infinite Arrangements. Version arXiv:2609.00603v1 (CC BY 4.0), distributed by arXiv under the Creative Commons Attribution 4.0 International licence, http://creativecommons.org/licenses/by/4.0/, as recorded in the arXiv licence metadata for this exact version and shown on the abstract page https://arxiv.org/abs/2609.00603v1. Full licence notice: \texttt{licenses/arxiv-2609.00603-CC-BY-4.0.txt}.\par\smallskip

\noindent\textit{Editorial scope.} Counted unit: the article's theorem thm:zonal-rigidity (source0.tex lines 1296--1333). The article's complete original proof of it is delivered with it (source0.tex lines 1335--1441, one \texttt{proof} environment of 107 lines). No mathematics is written, repaired or re-derived: the statement and the proof are the article's own lines, in the article's own notation, and no retained line is substituted or re-flowed.  Retained uncounted as the source-local context the proof needs: lines 788--793.  Nothing was omitted from any retained span, so the package carries no omission record.  No retained line contains a citation, so the wrapper carries no bibliography and no bibliography entry.  No retained line contains a figure environment, an image-inclusion command or drawing code, so no image is delivered and the package declares no asset dependency.  Editorial formatting only, in addition: an AMS \texttt{amsart} preamble that restates the article's own declarations and macros used by the retained lines, namely the environments \texttt{lemma}, \texttt{theorem} on separate counters, together with the standard packages those lines need.  The retained environments print in this document as Lemma 1, Theorem 1 in order of appearance, whereas the article prints the same items with the numbering of its own full document.  The complete native source of the article as distributed in the arXiv e-print for this exact version is bundled unchanged as \texttt{sources/209044/source0.tex}.
\begin{lemma}[Uniform orientation within the random plane]
\label{lem:uniform-frame-phase}
Let \(P=\operatorname{span}\{U,V\}\), equipped with the orientation induced
by the ordered frame.  Conditional on the oriented plane \(P\), the frame's
in-plane rotation is Haar-uniform on \(SO(2)\).
\end{lemma}

\begin{theorem}[Zonal magnitude rigidity]
\label{thm:zonal-rigidity}
Let \(p\geq3\), fix \(a\in\sphere^{p-1}\), and let
\(\psi:(-1,0)\cup(0,1)\to\R\) be Borel and finite almost everywhere.
Define the measurable tangent field
\begin{equation}
  T_\psi(u)
  :=
  \psi(a^\top u)\{a-(a^\top u)u\}
  \label{eq:zonal-field}
\end{equation}
away from \(a^\perp\), with an arbitrary definition on that null set.
The following are equivalent.
\begin{enumerate}[label=\textup{(\roman*)}]
\item For Haar-almost every oriented two-plane \(P\),
  \[
    \Law\{T_\psi(U)^\top V\mid P\}=\Cauchy(0,1).
  \]
\item
  \[
    |\psi(t)|=\frac1{|t|}
    \quad\text{for Lebesgue-almost every }t\in(-1,1).
  \]
\end{enumerate}
If \(T_\psi\) is also antipodally odd, then necessarily
\begin{equation}
  T_\psi(u)
  =
  \varepsilon(|a^\top u|)
  \left\{\frac{a}{a^\top u}-u\right\},
  \qquad
  \varepsilon:(0,1)\to\{-1,1\}
  \label{eq:zonal-sign-family}
\end{equation}
for an arbitrary measurable sign function, and every such field satisfies
\textup{(i)}.  If \(\psi\) is continuous on each half-interval, the sign is
constant and only the reciprocal field and its radial dual remain.
\end{theorem}

\begin{proof}
For a nonexceptional oriented plane \(P\), choose a positive orthonormal
basis so that
\[
  \operatorname{proj}_P(a)=r e_1,
  \qquad 0<r<1.
\]
Here \(r<1\) almost surely because \(p\geq3\), and the aligned bases can be
chosen measurably on a countable Grassmannian atlas.  Conditional uniformity
of the phase from \cref{lem:uniform-frame-phase} is therefore preserved.
With
\[
  U_\theta=e_1\cos\theta+e_2\sin\theta,
  \qquad
  V_\theta=-e_1\sin\theta+e_2\cos\theta,
\]
where \(\theta\) is uniform modulo \(2\pi\), direct calculation gives
\begin{equation}
  Z_r(\theta)
  :=
  T_\psi(U_\theta)^\top V_\theta
  =
  -r\sin\theta\,\psi(r\cos\theta).
  \label{eq:zonal-projection}
\end{equation}
For a Haar plane,
\(r^2\sim\Beta(1,(p-2)/2)\); its density is strictly positive throughout
\((0,1)\).  The conditional law in \cref{eq:zonal-projection} depends only
on \(r\), so \textup{(i)} is equivalent to \(Z_r\) being standard Cauchy
for Lebesgue-almost every \(r\in(0,1)\).

Fix \(0<s<1\).  Since \(\cos\theta\) has the arcsine density, the absolute
\(s\)-moment at every such \(r\) is
\begin{equation}
  \E|Z_r|^s
  =
  \frac1\pi
  \int_{-r}^r
  (r^2-t^2)^{(s-1)/2}|\psi(t)|^s\,\dd t.
  \label{eq:zonal-moment-transform}
\end{equation}
For a standard Cauchy variable this moment equals
\(\sec(\pi s/2)\).  The same beta integral shows that replacing
\(|\psi(t)|^s\) in \cref{eq:zonal-moment-transform} by \(|t|^{-s}\)
also gives \(\sec(\pi s/2)\).

For \(t>0\), put
\[
  d_s(t)
  :=
  \frac{|\psi(t)|^s+|\psi(-t)|^s}{2}-t^{-s}.
\]
It follows that
\begin{equation}
  \int_0^r
  (r^2-t^2)^{\beta-1}d_s(t)\,\dd t=0
  \quad\text{for almost every }r,
  \qquad
  \beta:=\frac{1+s}{2}\in(1/2,1).
  \label{eq:zonal-abel-zero}
\end{equation}
The moment identity at any good radius \(r_0\) shows that
\[
  \int_0^{r_0}
  \bigl\{|\psi(t)|^s+|\psi(-t)|^s\bigr\}\,\dd t<\infty:
\]
indeed, the kernel in \cref{eq:zonal-moment-transform} is bounded below by
\(r_0^{s-1}\) on \((0,r_0)\).  Since \(t^{-s}\) is integrable at zero,
\(d_s\) is locally integrable.  In
\cref{eq:zonal-abel-zero}, substitute \(x=r^2\), \(y=t^2\), and
\[
  f_s(y):=\frac{d_s(\sqrt y)}{2\sqrt y}.
\]
The resulting equation is the vanishing of the
Riemann--Liouville fractional integral \(I_{0+}^{\beta}f_s\).  Applying
\(I_{0+}^{1-\beta}\) and using the semigroup identity gives the next
display.  The interchange is justified first for \(|f_s|\) by Tonelli and
then for \(f_s\) by Fubini; the inner integral is the beta integral.  Thus
\[
  \int_0^x f_s(y)\,\dd y=0
  \quad\text{for almost every }x.
\]
The left side is absolutely continuous, so it vanishes everywhere and
\(d_s=0\) almost everywhere.

Two exponents are needed to identify two unknown magnitudes.  Use
\(s=1/4\) and \(s=1/2\) on their common full-measure set.  If
\[
  A=t|\psi(t)|,\qquad C=t|\psi(-t)|,
\]
then \(A^{1/4}+C^{1/4}=2\) and
\(A^{1/2}+C^{1/2}=2\).  Writing \(x=A^{1/4}\) and \(y=C^{1/4}\), these
relations give \(x+y=2\) and \(x^2+y^2=2\), hence the double solution
\(x=y=1\).  This proves \textup{(ii)}.

Conversely, under \textup{(ii)},
\(|Z_r(\theta)|=|\tan\theta|\) almost everywhere.  The transformation
\(\theta\mapsto-\theta\) preserves \(\cos\theta\) and reverses the sign in
\cref{eq:zonal-projection}; hence \(Z_r\) is symmetric and has the
half-Cauchy absolute-value law.  It is therefore standard Cauchy.

Finally, \(T_\psi(-u)=-T_\psi(u)\) is equivalent to
\(\psi(-t)=-\psi(t)\).  Combining this with \textup{(ii)} gives
\cref{eq:zonal-sign-family}.  Continuity makes the
\(\{-1,1\}\)-valued function \(\varepsilon\) constant on the connected
interval \((0,1)\).
\end{proof}

\end{document}
